\begin{ouverture}[equadiff]{6}{Primitives et équations différentielles}{Analyse et évolution}{Familles de solutions · Condition initiale · Euler}
\objectif{Reconnaître et déterminer des primitives.}
\objectif{Vérifier une solution d’équation différentielle.}
\objectif{Résoudre les équations $y\,\!^{\prime}=ay+b$.}
\objectif{Comparer modèle continu et approximation par Euler.}
\prerequis{Dérivation, exponentielle et récurrences affines.}
\end{ouverture}
\essentiel

\essSection{1}{Remonter d'une dérivée à une fonction}
\begin{definition}
Une primitive de $f$ sur un intervalle $I$ est une fonction $F$ dérivable telle que $F\,\!^{\prime}=f$ sur $I$.
Si $F$ est une primitive, toutes les primitives sur cet intervalle sont $F+C$, où $C$ est une constante réelle.
\end{definition}
\begin{demonstration}
Si $F$ et $G$ sont deux primitives, $(F-G)'=0$ sur l'intervalle. Une fonction de dérivée nulle sur un intervalle est constante, donc $F-G=C$.
L'hypothèse d'intervalle compte : sur deux morceaux disjoints, les constantes peuvent être différentes.
\end{demonstration}
\begin{center}\begin{tabular}{c|c}Fonction & Une primitive\\\hline
$x^n$ ($n\in\N$)&$x^{n+1}/(n+1)$\\
$1/x$ sur $]0,+\infty[$&$\ln x$\\
$\mathrm e^{ax+b}$ ($a\ne0$)&$\mathrm e^{ax+b}/a$\\
$2uu'$&$u^2$\\
$u'\mathrm e^u$&$\mathrm e^u$\\
$u'/u$ ($u>0$)&$\ln u$\end{tabular}\end{center}
\begin{exemple}[Reconnaître une forme]
Une primitive de $6x(x^2+1)^2$ est $(x^2+1)^3$, car sa dérivée vaut $3(x^2+1)^2\times2x$.
Pour $2x/(x^2+1)$, une primitive est $\ln(x^2+1)$.
Toujours dériver la proposition pour la contrôler.
\end{exemple}
\essSection{2}{Une équation portant sur une fonction}
Une équation différentielle relie une fonction inconnue et ses dérivées. Vérifier une solution consiste à calculer la dérivée puis à substituer les expressions dans l'équation, pour tout point du domaine.
Une condition initiale impose une valeur, par exemple $y(0)=y_0$.
\begin{theoreme}[$y\,\!^{\prime}=ay$]
Les solutions sur $\R$ sont $y(t)=C\mathrm e^{at}$ avec $C\in\R$. Avec $y(0)=y_0$, la solution unique est $y(t)=y_0\mathrm e^{at}$.
\end{theoreme}
\begin{demonstration}
Pour une solution $y$, poser $g(t)=y(t)\mathrm e^{-at}$. Alors $g\,\!^{\prime}(t)=(y\,\!^{\prime}-ay)\mathrm e^{-at}=0$.
Donc $g=C$, d'où $y=C\mathrm e^{at}$. La réciproque se vérifie en dérivant.
\end{demonstration}
\essSection{3}{Ajouter un terme constant}
\begin{savoirfaire}[$y\,\!^{\prime}=ay+b$]
Si $a\ne0$, la solution constante est $\ell=-b/a$. Poser $z=y-\ell$ : $z\,\!^{\prime}=az$.
Ainsi $y(t)=\ell+C\mathrm e^{at}$ ; avec $y(0)=y_0$, $C=y_0-\ell$.
Si $a=0$, l'équation est $y\,\!^{\prime}=b$, donc $y(t)=bt+C$.
\end{savoirfaire}
\begin{exemple}[Refroidissement]
Le modèle $T\,\!^{\prime}=-0{,}2(T-20)$, $T(0)=80$ donne $T(t)=20+60\mathrm e^{-0{,}2t}$.
La température tend vers $20$, celle de l'environnement supposée constante. Le coefficient $0{,}2$ dépend de l'unité de temps.
\end{exemple}
\figurelivre{c06-equadiff.pdf}{La solution exacte du refroidissement et les valeurs d'Euler avec pas $1$. Euler utilise des segments de tangente et donne une approximation, pas une identité.}
\essSection{4}{Du continu au discret : Euler}
Pour $y\,\!^{\prime}=f(t,y)$ et un pas $h>0$, la méthode d'Euler donne
$t_{n+1}=t_n+h$ et $y_{n+1}=y_n+h f(t_n,y_n)$.
Pour le refroidissement, $T_{n+1}=(1-0{,}2h)T_n+4h$ : on retrouve une suite arithmético-géométrique.
\begin{python}
def euler_refroidissement(h, n):
    t, temperature = 0.0, 80.0
    points = [(t, temperature)]
    for k in range(n):
        temperature += h*(-0.2*temperature + 4)
        t += h
        points.append((t, temperature))
    return points
\end{python}
Un pas trop grand peut produire un comportement artificiel. Réduire $h$ et comparer à une solution connue permet de contrôler l'approximation. Les équations non linéaires, comme la logistique, peuvent être explorées sur un exemple guidé sans théorie générale exigée.

\exercices
\markright{Prendre confiance}
\begingroup\renewcommand{\otBesoinPlace}[1]{}\begin{minipage}{\linewidth}
\rubrique{Prendre confiance}
\begin{exercice}[1]{Primitives usuelles}

\begin{questions}
\item Donner toutes les primitives de $3x^2-4$.
\item Donner celles de $\mathrm e^{2x}$.
\item Donner celles de $1/x$ sur $]0,+\infty[$.
\end{questions}
\end{exercice}
\begin{exercice}[1]{Reconnaître une composée}

\begin{questions}
\item Donner une primitive de $2x\mathrm e^{x^2}$.
\item Donner une primitive de $2x/(x^2+3)$.
\item Donner une primitive de $4x(x^2+1)$.
\end{questions}
\end{exercice}
\begin{exercice}[1]{Condition initiale}

\begin{questions}
\item Trouver $F$ telle que $F\,\!^{\prime}(x)=2x$ et $F(0)=3$.
\item Trouver $G$ telle que $G\,\!^{\prime}(x)=\mathrm e^x$ et $G(1)=0$.
\item Pourquoi la constante ne doit-elle pas être oubliée ?
\end{questions}
\end{exercice}
\begin{exercice}[1]{Vérifier une solution}
$y(t)=4+3\mathrm e^{-2t}$.
\begin{questions}
\item Calculer $y\,\!^{\prime}$.
\item Vérifier $y\,\!^{\prime}=-2y+8$.
\item Donner la condition initiale satisfaite.
\end{questions}
\end{exercice}
\end{minipage}\endgroup
\clearpage
\markright{Appliquer et relier}
\begingroup\renewcommand{\otBesoinPlace}[1]{}\begin{minipage}{\linewidth}
\rubrique{Appliquer et relier}
\begin{exercice}[2]{Résoudre sans terme constant}

\begin{questions}
\item Résoudre $y\,\!^{\prime}=3y$, $y(0)=2$.
\item Résoudre $z\,\!^{\prime}=-z/2$, $z(0)=10$.
\item Donner leurs limites quand $t\to+\infty$.
\end{questions}
\end{exercice}
\begin{exercice}[2]{Déplacer un équilibre continu}
$y\,\!^{\prime}=-0{,}5y+6$, $y(0)=2$.
\begin{questions}
\item Trouver une solution constante.
\item Donner la solution du problème.
\item Déterminer le temps pour dépasser $11$.
\end{questions}
\end{exercice}
\begin{exercice}[2]{Cas $a=0$}
On considère $y\,\!^{\prime}=ay+5$ avec $a=0$ et $y(2)=9$.
\begin{questions}
\item Pourquoi la formule avec $-b/a$ ne s’applique-t-elle pas ?
\item Résoudre directement.
\item Vérifier la condition donnée.
\end{questions}
\end{exercice}
\begin{exercice}[2]{Euler à la main}
$y\,\!^{\prime}=-y$, $y(0)=1$. On choisit $h=0{,}5$.
\begin{questions}
\item Écrire la récurrence d’Euler.
\item Calculer les valeurs aux temps $0{,}5$ et $1$.
\item Comparer à la solution exacte au temps $1$.
\end{questions}
\end{exercice}
\end{minipage}\endgroup
\clearpage
\markright{Chercher et argumenter}
\begingroup\renewcommand{\otBesoinPlace}[1]{}\begin{minipage}{\linewidth}
\rubrique{Chercher et argumenter}
\begin{exercice}[3]{Pourquoi une seule constante ?}
$F$ et $G$ sont primitives de la même fonction sur un intervalle $I$.
\begin{questions}
\item Dériver $F-G$.
\item En déduire la forme de $G$.
\item Donner un exemple montrant que l’hypothèse d’intervalle est utile.
\end{questions}
\end{exercice}
\begin{exercice}[3]{Pas trop grand}
Pour $y\,\!^{\prime}=-2y$, $y(0)=1$, Euler donne $y_{n+1}=(1-2h)y_n$.
\begin{questions}
\item Étudier les valeurs obtenues pour $h=1$.
\item Comparer au signe de la solution exacte.
\item Que se passe-t-il pour $h=0{,}25$ ?
\end{questions}
\end{exercice}
\begin{exercice}[3]{Logistique : vérification guidée}
$y(t)=100/(1+4\mathrm e^{-t})$.
\begin{questions}
\item Calculer $y(0)$ et sa limite en $+\infty$.
\item Vérifier $y\,\!^{\prime}=y(1-y/100)$.
\item Pourquoi cette vérification ne constitue-t-elle pas une méthode générale de résolution ?
\end{questions}
\end{exercice}
\begin{exercice}[3]{Modèles discret et continu}
Une quantité suit $u_n=50+30(0{,}8)^n$.
\begin{questions}
\item Trouver $k$ pour que $y(t)=50+30\mathrm e^{-kt}$ coïncide avec $u_n$ aux temps entiers.
\item Écrire une équation différentielle satisfaite par $y$.
\item Ce modèle continu est-il l’Euler de pas $1$ de cette équation ?
\end{questions}
\end{exercice}
\end{minipage}\endgroup
